\section{Reconstruction results}\label{sec:results}

\begin{table*}
\caption{Reconstruction benchmark at $90\%$ missing data (mean over two random masks; 1500 test entries and 64 samples per entry). The correlation $\rho$ and the NRMSE refer to the posterior median; CRPS, Eq.~\eqref{eq:crps}, is normalized by the standard deviation of the target; ``cov.'' is the empirical coverage of the central $90\%$ interval; ``sign'' is the fraction of samples with the correct sign of the first target ($\cos\theta$ in C1, $\sin\theta$ in C2 and C3); $W_1$ is the sliced Wasserstein distance between imputed and true Poincar\'e sections (C1 only). In C1, $\rho$, NRMSE, CRPS, and cov.\ are averaged over the two targets $(c,\dot\theta)$. The original FPRM is the reimplemented GPR on the delay window of the complete channel (Sec.~\ref{sec:tasks}); in C3 with phase, the phase is appended to its input. In C2 the original FPRM and GPR (full context) converge to the same noise-dominated predictor, hence their identical rows. Results for every missing fraction and random mask are listed in Tables~S1--S4 of the supplementary material.}
\label{tab:fprm}
\PIARlayout{\renewcommand{\baselinestretch}{1}\footnotesize}{}%
\begin{ruledtabular}
\begin{tabular}{lcccccc}
Imputer & $\rho$ & NRMSE & CRPS & cov. & sign & $W_1$\\
\hline
\multicolumn{7}{l}{\textit{C1: $s\to(c,\dot\theta)$, forced}}\\
static GPR & 0.116 & 0.996 & 0.570 & 0.91 & 0.549 & 0.524\\
original FPRM (GPR) & 0.893 & 0.449 & 0.167 & 0.90 & 0.893 & 0.079\\
GPR (full context) & 0.971 & 0.216 & 0.068 & 0.90 & 0.971 & 0.044\\
cVAE & 0.917 & 0.386 & 0.172 & 0.86 & 0.815 & 0.066\\
DDPM & 0.990 & 0.135 & 0.042 & 0.93 & 0.965 & 0.024\\
NSF & 0.980 & 0.189 & 0.054 & 0.80 & 0.953 & 0.020\\
NSF + physics & 0.981 & 0.171 & 0.037 & 0.66 & 0.973 & 0.019\\
\multicolumn{7}{l}{\textit{C2: $c\to s$, Langevin}}\\
static GPR & 0.054 & 1.016 & 0.581 & 0.90 & 0.505 & --\\
original FPRM (GPR) & 0.031 & 1.008 & 0.582 & 0.91 & 0.502 & --\\
GPR (full context) & 0.031 & 1.008 & 0.582 & 0.91 & 0.502 & --\\
cVAE & 0.635 & 0.777 & 0.378 & 0.85 & 0.659 & --\\
DDPM & 0.517 & 0.961 & 0.270 & 0.91 & 0.674 & --\\
NSF & 0.488 & 0.986 & 0.287 & 0.78 & 0.672 & --\\
NSF + physics & 0.032 & 1.389 & 0.735 & 0.59 & 0.527 & --\\
\multicolumn{7}{l}{\textit{C3: $c\to s$, forced, with phase}}\\
original FPRM (GPR) + phase & 0.989 & 0.145 & 0.047 & 0.95 & 0.972 & --\\
NSF & 0.975 & 0.224 & 0.067 & 0.79 & 0.941 & --\\
NSF + physics & 0.488 & 0.879 & 0.491 & 0.35 & 0.715 & --\\
\multicolumn{7}{l}{\textit{C3: $c\to s$, forced, no phase}}\\
original FPRM (GPR) & $-0.025$ & 1.020 & 0.587 & 0.99 & 0.500 & --\\
NSF & 0.427 & 1.045 & 0.329 & 0.84 & 0.640 & --\\
NSF + physics & 0.351 & 1.206 & 0.664 & 0.38 & 0.599 & --\\
\end{tabular}

\end{ruledtabular}
\end{table*}

\begin{figure*}
\includegraphics[width=\textwidth]{figures/fig_fprm.pdf}
\caption{Reconstruction scores versus missing fraction. (a) NSF and (b) GPR (full context): posterior samples of the spring force $s$ versus the complete collar height $c$ for the missing entries of C2 at $90\%$ missing data that have no observed neighbor (first random mask; logarithmic color scale); the flow places its samples on both branches $s=\pm(1-c^{2})^{1/2}$ (black), whereas the GPR averages them. (c) CRPS versus the missing fraction for C2. (d) CRPS and (e) sliced Wasserstein distance between imputed and true Poincar\'e sections for C1. (f) CRPS for C3 with (filled markers) and without (open markers, dotted lines) the forcing phase; the original FPRM receives the phase in the first variant. Markers show the mean over two random masks and bars the half-range; the CRPS is normalized by the standard deviation of the target.}
\label{fig:fprm}
\end{figure*}

Table~\ref{tab:fprm} lists all metrics at $90\%$ missing data, the most demanding setting, and Fig.~\ref{fig:fprm} shows how the scores depend on the missing fraction. The results are presented in the order dictated by the identifiability analysis: a trivial control, an identifiable task, a task made ambiguous by symmetry, and its driven analogue, in which the forcing phase can lift the ambiguity.

\subsection{A trivial control (C0)}\label{sec:res_c0}

With the angle complete, a Savitzky--Golay derivative (window of 11 samples, degree 4) reconstructs the angular velocity with $\rho=0.9997$ and $\mathrm{NRMSE}=0.023$--$0.025$ at every missing fraction, since it never uses the incomplete channel. The flow, trained only on the observed pairs and conditioned on the full vector of Eq.~\eqref{eq:context}, is less accurate: $\rho=0.998$ at $30\%$ and $0.65$ at $90\%$ missing data. Its input is the unwrapped angle, whose secular drift during rotations makes it a poor observable for delay reconstruction and inflates the selected delay ($\tau=26$ samples). Reconstruction from the angle is therefore not a meaningful benchmark, and the tasks below use observables that are not related by differentiation.

\subsection{Identifiable chaotic dynamics (C1)}\label{sec:res_c1}

When the spring force is complete, the magnitude $|c|=(1-s^{2})^{1/2}$ follows from the holonomic identity, and the informative targets are the branch, i.e., the sign of $\cos\theta$, and the angular velocity. Every imputer that receives the full conditioning vector reconstructs $(c,\dot\theta)$ accurately at $30\%$ missing data ($\rho\ge0.994$), and the differences emerge as data are removed [Fig.~\ref{fig:fprm}(d)]. The original FPRM, which relies on the delay window alone, remains at $\rho\approx0.87$--$0.89$ and $\mathrm{NRMSE}\approx0.45$--$0.49$ at all fractions.

The accuracy of the original FPRM is limited not by identifiability, which holds for this task [Fig.~\ref{fig:fprm_scheme}(b)], but by the number of training points. A nearest-neighbor predictor on the same delay window reaches $\rho=0.98$ at $30\%$ missing data when all $2.8\times10^{4}$ observed entries serve as its library, but only $\rho=0.84$ when the library is restricted to 1200 points, the size of the training set of the Gaussian process. Resolving a folded chaotic attractor from a delay window alone requires a dense library, as in convergent cross mapping~\cite{Sugihara2012}. Fitting the same Gaussian process to all labeled samples at $90\%$ missing data (about 3900--4000 per mask) confirms this diagnosis: its correlation rises to $\rho=0.936$ and its CRPS falls from $0.167$ to $0.112$, with a coverage of $0.93$.

Among the purely data-driven engines that receive the full conditioning vector, the flow and the cVAE are on par at $30\%$ (CRPS $0.012$ and $0.013$), the flow is best at $60\%$ ($0.021$), and at $90\%$ the diffusion model is better ($0.042$ against $0.054$), more accurate in the median (NRMSE $0.135$ against $0.189$), and better calibrated (coverage $0.93$ against $0.80$). The physics-regularized flow attains the lowest CRPS at every fraction ($0.011$, $0.018$, and $0.037$, that is, $13\%$, $11\%$, and $31\%$ below the unregularized flow), and at $90\%$ it reduces the fraction of samples on the wrong branch from $4.7\%$ to $2.7\%$. Against the original FPRM, its CRPS is $4.5$ times lower at $90\%$ missing data ($3.6$ and $5.7$ times for the two masks), and still three times lower ($2.5$ and $3.7$ times) when the Gaussian process is trained on every label. The GPR that receives the same full conditioning vector closes much of this gap (CRPS $0.068$), which shows that a large part of the gain comes from the conditioning information, i.e., from the observed neighbors and the forcing phase, rather than from the density model alone. The price of the physics term is calibration: the coverage of the nominal $90\%$ intervals falls to $0.66$, so that the regularized flow is overconfident.

The benefit becomes tangible when the dynamics itself is reconstructed. In Fig.~\ref{fig:reconstruction}, every missing entry of an eight-period window at $90\%$ missing data, 236 entries of each target, is imputed with the regularized flow and with the original FPRM. The posterior medians of the flow follow the continuous trajectory through rapid rotations and slow librations, with root-mean-square errors of $0.22$ in $c$ and $0.17$ in $\dot\theta\sqrt{L/g}$, against $0.34$ and $0.45$ for the original FPRM, whose Gaussian means scatter between the branches. Imputing $(c,\dot\theta)$ at all 1137 missing section points and combining them with the noise-free spring force rebuilds the Poincar\'e section. The flow recovers the folded, fractal structure of the attractor ($W_1=0.023$), whereas the original FPRM smears it into a diffuse cloud ($W_1=0.063$). On the benchmark, whose section is taken one sample later, at $\Omega t=2\pi/32$, because its sampled record starts one sample after a forcing period, both flows produce the sections closest to the true one ($W_1=0.019$--$0.020$ at $90\%$, against $0.024$ for the diffusion model and $0.079$ for the original FPRM) [Fig.~\ref{fig:fprm}(e)].

\begin{figure*}
\includegraphics[width=\textwidth]{figures/fig_reconstruction.pdf}
\caption{Chaotic dynamics recovered from $10\%$ of the collar-height and angular-velocity records (task C1, first random mask). (a) Collar height and (b) angular velocity over eight forcing periods: continuous true trajectory (gray line), observed samples (black dots), posterior median and central $90\%$ interval of the physics-regularized flow at every missing sample (blue), and Gaussian mean of the original FPRM (orange squares). (c) Physics-regularized flow and (d) original FPRM: stroboscopic Poincar\'e section at the phase $\Omega t=0$ of the drive, rebuilt at the 1137 missing section points from the noise-free spring force and one imputation of $(c,\dot\theta)$ (colored), over the true section (gray); $W_1$ is the sliced Wasserstein distance in $(s,c,\dot\theta)$.}
\label{fig:reconstruction}
\end{figure*}

\subsection{Symmetry-induced ambiguity (C2)}\label{sec:res_c2}

When the collar height is complete, the original FPRM returns $\rho\approx0$ at every missing fraction, as predicted by the reflection symmetry, and its CRPS does not improve as data are added [Fig.~\ref{fig:fprm}(c)]. At $30\%$ missing data the observed neighbors resolve the sign for almost every entry, all imputers that receive the full conditioning vector succeed ($\rho\ge0.994$), and the two flows have the lowest errors (NRMSE $0.074$; CRPS $0.023$ without and $0.020$ with the physics term). At $90\%$, about half of the missing entries ($0.9^{6}\approx53\%$ expected) have no observed neighbor within three samples, and their exact conditional distribution is bimodal. The flow represents both branches [Fig.~\ref{fig:fprm}(a)], whereas the GPR with the full conditioning vector, whose marginal-likelihood optimization (with a single start) converges to a noise-dominated solution, returns a broad Gaussian centered at zero [Fig.~\ref{fig:fprm}(b)]. The CRPS values of the flow and of the diffusion model ($0.29$ and $0.27$) are about half that of the GPR ($0.58$). In this regime the correlation of the posterior median is a misleading score, because the median of a nearly symmetric bimodal distribution lies between the branches: the cVAE, whose Gaussian decoder tends to merge the branches, reaches a higher correlation ($0.63$) but a worse CRPS ($0.38$).

\subsection{The forcing phase restores identifiability (C3)}\label{sec:res_c3}

In this task the original FPRM, with and without the phase appended to its delay window, is compared with the two flows. In the driven system the phase lifts the ambiguity [Fig.~\ref{fig:fprm}(f)], as anticipated by the cross-mapping test of Fig.~\ref{fig:fprm_scheme}(d). With the phase, the original FPRM reconstructs $s$ with $\rho=0.989$ and CRPS $0.047$ at $90\%$ missing data, better than the flow ($\rho=0.975$, CRPS $0.067$), whose conditional density is harder to estimate from about 4000 labeled samples; at $30\%$ and $60\%$ the flow is more accurate (NRMSE $0.069$ and $0.090$, against $0.163$ and $0.160$). Without the phase, the original FPRM fails completely ($\rho\approx0$), and the flow, which then relies on the observed neighbors alone, falls from $\rho=0.993$ at $30\%$ to $0.43$ at $90\%$.

\subsection{When the physics helps and when it hurts}\label{sec:res_physics}

At $30\%$ missing data the physics residuals lower the CRPS of the flow in C1, C2, and C3 with the phase (by $13\%$, $14\%$, and $21\%$, respectively) and change it by less than $2\%$ in C3 without the phase; at this fraction the observed neighbors resolve the branch of most entries, and the residuals act only as a mild regularizer of the density. At larger missing fractions they continue to help in C1, where the gravitational and forcing terms of Eq.~\eqref{eq:res_eom} discriminate the sign of $\cos\theta$ and the kinematic relation fixes the angular velocity, but they are harmful in C2 and C3. At $90\%$ missing data the regularized flow attains only $\rho=0.03$ in C2 and $0.35$ in C3 without the phase, against $0.49$ and $0.43$ for the unregularized flow, and its coverage drops to $0.59$ and $0.38$; in C2 one mask collapses already at $60\%$ ($\rho=-0.04$). The same failure appears in C3 with the phase, where the target is identifiable: at $60$--$90\%$ the CRPS is worse than without the physics term in three of the four runs, and one mask collapses at $90\%$ ($\rho=0.05$, against $0.98$ without the physics term).

The pattern follows from the structure of the residuals. Where the target is ambiguous, the residuals are nearly invariant under $\hat s\to-\hat s$ and cannot select a branch, but they penalize the probability that a flow with a unimodal base must place between the two branches~\cite{Cornish2020}. The flow can lower this penalty by collapsing onto one branch, after which it assigns high confidence to a single, often wrong, sign. In C3, moreover, the only sign-discriminating term of Eq.~\eqref{eq:res_c3}, $f(t)\,\hat s^{3}$, competes with the noise that the numerical second derivative contributes to the residual, $\hat s^{2}\,\delta\ddot c$, where $\delta\ddot c\approx0.107$ is the standard deviation of the Savitzky--Golay estimate of $\ddot c$ for a measurement noise of $0.01$. As an order-of-magnitude estimate, the discriminating term falls below this noise whenever $A|\hat s\cos\Omega t|<\delta\ddot c$, i.e., $|\hat s\cos\Omega t|<0.13$, which is the case for $24\%$ of the samples on the attractor; the residual is therefore poorly conditioned even when the phase is available. The collapse also escapes early stopping, because the validation likelihood is computed on labeled entries and does not see the imputations of unlabeled ones. These observations are related to the known failure modes of physics-informed training, in which soft constraints make the loss landscape ill conditioned~\cite{Krishnapriyan2021,WangTengPerdikaris2021}; here, in addition, the residual carries no information about the branch.

\subsection{Flow versus diffusion model}\label{sec:res_engine}

The diffusion model is the strongest competitor of the flow, and the choice between the two engines rests on the following evidence. Without physical residuals neither engine dominates: the flow attains a lower CRPS than the DDPM at $30\%$ and $60\%$ missing data in task C1 (by $18\%$) and at $30\%$ in task C2 (by $22\%$), the DDPM attains a lower CRPS at $90\%$ (by $22\%$ in C1 and $6\%$ in C2), and the two agree to within $3\%$ at $60\%$ in C2; each strict ranking holds for both masks. The DDPM reached its lowest validation loss within the last $10\%$ of its 6000 training steps in every run, so that its scores might still improve slightly with a longer budget. The intervals of the DDPM are conservative (coverage $0.93$--$0.97$ in C1), whereas those of the flow become too narrow as data are removed ($0.80$ at $90\%$); the flow, however, reproduces the invariant measure more faithfully at every missing fraction, with sliced distances on the Poincar\'e section smaller by $16$--$24\%$.

The decisive differences are structural. With the architectures of Table~\ref{tab:hyper} on one processor thread, drawing 64 samples for 1500 missing entries of task C1 took $1.8$~s with the flow and $47$~s with the DDPM, because the flow evaluates its conditioner six times per sample (two autoregressive passes in each of three transforms), whereas the DDPM evaluates its denoiser once per noise level, i.e., 100 times. For the scalar target of task C2, whose spline parameters depend only on the conditioning vector and are computed once per entry, the ratio reaches three orders of magnitude ($0.02$~s against $27$~s). Accelerated samplers reduce the number of denoising steps at some cost in sample quality~\cite{SongMengErmon2021}, which narrows this gap without closing it.

A physics-regularized training step needs gradients through the samples: the flow obtains them from one reparameterized pass, whereas the DDPM requires backpropagation through its entire reverse chain, which multiplied the cost of differentiating a penalty evaluated on the samples by more than $30$ in task C1 and by more than $200$ in task C2 (Table~S8 of the supplementary material). Physics-informed diffusion models avoid this cost by evaluating the residuals on the denoised mean estimate at each training noise level, or on a few accelerated sampling steps, rather than on full-chain samples of the model~\cite{Bastek2025}. Moreover, the exact likelihood of a diffusion model requires the integration of a probability-flow ordinary differential equation~\cite{Song2021}, whereas the flow returns it directly. Residuals evaluated on samples of the model, as in Eq.~\eqref{eq:loss}, were therefore affordable only for the flow, and the physics-regularized flow attains the lowest mean CRPS of all imputers in task C1 at every missing fraction: $27$--$29\%$ below that of the DDPM at $30\%$ and $60\%$ missing data ($25$--$31\%$ for the individual masks), and $11\%$ below at $90\%$, where the two masks disagree (CRPS $0.0454$ and $0.0295$, against $0.0447$ and $0.0395$ for the DDPM). The price is calibration: every flow in this study covered less than the nominal $90\%$, the regularized ones even less, a deficit that must be removed, for instance by recalibration on held-out entries~\cite{Kuleshov2018}, before the intervals are used.
